\documentclass{report}

\input{~/dev/latex/template/preamble.tex}
\input{~/dev/latex/template/macros.tex}

\title{\Huge{}}
\author{\huge{Nathan Warner}}
\date{\huge{}}
\fancyhf{}
\rhead{}
\fancyhead[R]{\itshape Warner} % Left header: Section name
\fancyhead[L]{\itshape\leftmark}  % Right header: Page number
\cfoot{\thepage}
\renewcommand{\headrulewidth}{0pt} % Optional: Removes the header line
%\pagestyle{fancy}
%\fancyhf{}
%\lhead{Warner \thepage}
%\rhead{}
% \lhead{\leftmark}
%\cfoot{\thepage}
%\setborder
% \usepackage[default]{sourcecodepro}
% \usepackage[T1]{fontenc}

% Change the title
\hypersetup{
    pdftitle={Embedded insights}
}

\begin{document}
    % \maketitle
        \begin{titlepage}
       \begin{center}
           \vspace*{1cm}
    
           \textbf{Embedded insights}
    
           \vspace{0.5cm}
            
                
           \vspace{1.5cm}
    
           \textbf{Nathan Warner}
    
           \vfill
                
                
           \vspace{0.8cm}
         
           \includegraphics[width=0.4\textwidth]{~/niu/seal.png}
                
           Computer Science \\
           Northern Illinois University\\
           United States\\
           
                
       \end{center}
    \end{titlepage}
    \tableofcontents
    \pagebreak 
    \unsect{Concepts}
    \begin{itemize}
        \item \textbf{Electronics}
            \begin{itemize}
                \item Voltage, current, resistance, and power
                \item Ohm’s law
                \item Digital high and low voltage levels
                \item GPIO input versus output
                \item Pull-up and pull-down resistors
                \item LEDs and current-limiting resistors
                \item Push buttons and switch bouncing
                \item Capacitors and basic RC circuits
                \item Analog versus digital signals
                \item Reading schematics and pinout diagrams
                \item Common ground and safe voltage limits
                \item Basic use of a multimeter and logic analyzer
            \end{itemize}
        \item \textbf{MCU Architecture}: Learn how a microcontroller differs from a desktop processor: 
            \begin{itemize}
                \item CPU core, flash, SRAM, and peripherals
                \item Memory-mapped I/O
                \item Registers and special-function registers
                \item Cortex-M4 register set
                \item Vector table
                \item Main stack pointer
                \item Reset sequence
                \item Privileged and unprivileged execution
                \item Exceptions and interrupts
                \item NVIC interrupt controller
                \item SysTick timer
                \item Memory protection unit
                \item Floating-point unit
                \item Little-endian representation
                \item Thumb and Thumb-2 instructions
            \end{itemize}
        \item \textbf{The embedded build process}
            \begin{itemize}
                \item Cross-compilation
                \item Preprocessor, compiler, assembler, and linker stages
                \item Object files and ELF files
                \item Sections such as .text, .rodata, .data, and .bss
                \item Symbols and relocation
                \item Linker scripts
                \item Startup assembly
                \item Vector-table placement
                \item Copying .data from flash to SRAM
                \item Zero-initializing .bss
                \item Map files
                \item Converting ELF to BIN or Intel HEX
                \item Flashing versus debugging
            \end{itemize}
        \item \textbf{Essential STM32 peripherals}: 
            \begin{enumerate}
                \item GPIO
                \item Clock system and RCC
                \item External interrupts
                \item General-purpose timers
                \item PWM
                \item UART
                \item ADC
                \item SPI
                \item I²C
                \item DMA
                \item Watchdogs
                \item RTC and low-power modes
            \end{enumerate}
            For every peripheral, understand:
            \begin{itemize}
                \item Which clock enables it
                \item Which GPIO alternate function it uses
                \item Which registers configure it
                \item How status flags work
                \item How polling differs from interrupts
                \item How DMA can move data without continuous CPU involvement
                \item What errors and edge cases can occur
            \end{itemize}
            Board exercises include LED blinking with a timer, button interrupts, UART command input, potentiometer sampling through the ADC, PWM brightness control, and reading a sensor through $\text{I}^{2}\text{C}$ or SPI.
        \item \textbf{Interrupts and concurrency}: 
            \begin{itemize}
                \item Interrupt service routines
                \item Interrupt latency
                \item Interrupt priorities
                \item Shared state between main code and an ISR
                \item Atomic operations and critical sections
                \item Race conditions
                \item Reentrancy
                \item Keeping ISRs short
                \item Deferred processing
                \item Ring buffers
                \item Producer-consumer design
                \item Why blocking inside an ISR is dangerous
                \item Why volatile alone does not solve concurrency
            \end{itemize}
            Board exercise: receive UART bytes in an interrupt, place them into a ring buffer, and process them from the main loop.
        \item \textbf{Time and real-time behavior}: 
            \begin{itemize}
                \item Busy waiting versus hardware timers
                \item Blocking versus nonblocking code
                \item Deadlines, latency, jitter, and throughput
                \item Periodic and event-driven tasks
                \item Finite-state machines
                \item Cooperative scheduling
                \item Monotonic timestamps
                \item Timer overflow
                \item Input capture and output compare
                \item Debouncing without blocking delays
            \end{itemize}
            Board exercise: run LED control, button handling, UART processing, and sensor sampling concurrently without calling a long blocking delay.
        \item \textbf{Hardware communication protocols}: 
            \begin{itemize}
                \item \textbf{UART}	Baud rate, frames, parity, buffering, overruns
                \item \textbf{SPI}	Controller/peripheral roles, chip select, CPOL, CPHA, full duplex
                \item \textbf{I²C}	Addresses, START/STOP, ACK/NACK, pull-ups, clock stretching
                \item \textbf{CAN}	Frames, arbitration, identifiers, error handling
                \item \textbf{USB}	Endpoints, enumeration, descriptors; learn later
            \end{itemize}
        \item \textbf{Drivers and firmware architecture}: Learn to separate:
            \begin{itemize}
                \item Application logic
                \item Hardware-independent modules
                \item Device drivers
                \item Board support code
                \item MCU peripheral access
                \item Vendor HAL or low-level libraries
            \end{itemize}
            Also understand:
            \begin{itemize}
                \item Hardware abstraction layers
                \item Dependency injection through function pointers
                \item State machines
                \item Event queues
                \item Error reporting
                \item Timeouts
                \item Configuration versus runtime state
                \item Interfaces suitable for host-based testing
            \end{itemize}
        \item \textbf{RTOS concept}: Do not begin with an RTOS. First become comfortable with interrupts, timers, state machines, and a superloop. Then learn
            \begin{itemize}
                \item Tasks and context switching
                \item Preemptive scheduling
                \item Priorities
                \item Queues
                \item Semaphores
                \item Mutexes
                \item Software timers
                \item Event groups
                \item Priority inversion
                \item Deadlock
                \item ISR-safe APIs
                \item Per-task stack sizing
                \item Heap-allocation strategies
                \item Tick frequency and tickless operation
            \end{itemize}
        \item \textbf{For employment, you should at least recognize:}: 
            \begin{itemize}
                \item Watchdog design
                \item Brownout and reset handling
                \item Defensive timeouts
                \item Fault recovery
                \item Firmware versioning
                \item Bootloaders and firmware updates
                \item CRCs and data integrity
                \item Power loss during flash writes
                \item Static analysis
                \item Coding standards such as MISRA C
                \item Unit, integration, and hardware-in-the-loop testing
                \item Resource limits
                \item Low-power design
                \item Security fundamentals
            \end{itemize}
        \item \textbf{Progression}: 
            \begin{center}
                \begin{tabularx}{\textwidth}{@{}lXX@{}}
                    \toprule
                    \textbf{Stage}&	\textbf{Board work}&	\textbf{Main concepts} \\
                    \midrule
                    1	&LED and button	&GPIO, C bit operations, schematics\\[2ex]
                    2	&Build outside the IDE	&GCC, ELF, linker script, startup code\\[2ex]
                    3	&Timer-driven LED	&clocks, timers, interrupts, NVIC\\[2ex]
                    4	&UART command console	&serial communication, ring buffers\\[2ex]
                    5	&ADC and PWM	&analog input, sampling, timer output\\[2ex]
                    6	&I²C/SPI &sensor	protocols, datasheets, reusable drivers\\[2ex]
                    7	&DMA data acquisition	&concurrency, buffering, throughput\\[2ex]
                    8	&Cooperative application&	state machines, events, nonblocking design\\[2ex]
                    9	&FreeRTOS application	&tasks, queues, synchronization\\[2ex]
                    10	&Portfolio project	&architecture, testing, fault handling \\
                    \bottomrule
                \end{tabularx}
            \end{center}
        \item \textbf{For job}
            \begin{itemize}
                \item Write and debug C confidently.
                \item Explain startup, linking, flash, SRAM, and the vector table.
                \item Configure basic peripherals from their documentation.
                \item Write interrupt-driven, nonblocking firmware.
                \item Use UART, SPI, I²C, timers, ADC, PWM, and DMA.
                \item Diagnose faults using GDB and hardware tools.
                \item Explain races, critical sections, queues, and RTOS scheduling.
                \item Structure firmware into testable drivers and application modules.
                \item Present one substantial board project with documentation and tests.
            \end{itemize}
    \end{itemize}

    \pagebreak 
    \subsection{Electronics}
    \begin{itemize}
        \item \textbf{GPIO}: GPIO stands for General-Purpose Input/Output. GPIO allows software running on a microcontroller to interact with electrical signals through its physical pins.
            \bigbreak \noindent 
            A GPIO pin can usually operate as:
            \begin{itemize}
                \item A digital input
                \item A digital output
                \item A connection to an internal peripheral, such as UART, SPI, I$^2$C, or a timer
                \item An analog input or output, when supported
            \end{itemize}
            Software reads and writes registers, the GPIO hardware translates those register operations into electrical behavior at a pin.
            \bigbreak \noindent 
            When configured as an output, the MCU drives the pin to one of two logic levels:
            \begin{center}
                \begin{tabularx}{\textwidth}{@{}XXX@{}}
                    \toprule
                    \textbf{Software value}&	\textbf{Logical state}&	\textbf{Approximate voltage} \\
                    \midrule
                    0	&Low	&0 V \\[2ex]
                    1	&High	&Supply voltage, often 3.3 V \\
                    \bottomrule
                \end{tabularx}
            \end{center}
            Writing a 1 to the pin drives it high, causing current to flow through the resistor and LED.
            \begin{center}
                GPIO pin -- resistor -- LED -- ground
            \end{center}
        \item \textbf{Active-high versus active-low}: The previous circuit is active-high because setting the GPIO high turns on the LED.
            \begin{center}
                3.3 V -- resistor -- LED -- GPIO pin
            \end{center}
            This is active-low. You must inspect the board schematic to determine which arrangement is used.
        \item \textbf{Push buttons}: A push button is simply a mechanical switch. It does not produce a voltage or send data by itself. It either:
            \begin{itemize}
                \item Connects two electrical points when pressed, or
                \item Disconnects two electrical points when released.
            \end{itemize}
            The core difference between a normally open (NO) and normally closed (NC) button is their default state when you are not touching them
            \begin{itemize}
                \item \textbf{NO}: The internal contacts are apart. The circuit is open. When pressed the contacts touch and close the circuit. Current flows to turn something on.
                \item \textbf{NC}: The internal contacts touch each other. The circuit is complete. When pressed the contacts pull apart and open the circuit. Power stops.
            \end{itemize}
            A \textbf{momentary circuit button} (usually called a momentary push-button switch) is an electrical switch that turns a circuit on or off only while you are actively pressing it.
            \bigbreak \noindent 
            The opposite of a momentary switch is a latching switch (also called a maintained switch).
            \bigbreak \noindent 
            The blue user button on your NUCLEO-F446RE is a momentary, normally-open switch:
        \item \textbf{GPIO as an input}: When configured as an input, the MCU does not intentionally drive the pin. Instead, it measures whether the external pin voltage represents a logical low or high.
            \bigbreak \noindent 
            For a 3.3 V MCU, the interpretation is approximately:
            \begin{itemize}
                \item \textbf{Near 0 V:}  logic 0
                \item \textbf{Near 3.3 V:}  logic 1
            \end{itemize}
            The exact voltage thresholds are given in the MCU datasheet.
        \item \textbf{Pull-up and pull-down resistors}: Pull-up and pull-down resistors are small components used in digital circuits to give a pin a defined default voltage level and stop it from "floating"
            \bigbreak \noindent 
            A floating pin is unconnected. It acts like an antenna, picking up stray electrical noise that causes random, unpredictable high or low readings
            \begin{itemize}
                \item \textbf{Pull up resistor}: A pull-up resistor connects the input pin directly to the positive voltage supply (\(V_{CC}\)). It holds the pin at a logical HIGH (1) when a switch or button is open.
                    \bigbreak \noindent 
                    When you close the switch, it connects the pin directly to ground. This overrides the weak pull-up resistor and pulls the reading to a logical LOW (0).
                    \begin{verbatim}
                        3.3 V
                        |
                        resistor
                        |
                        GPIO input -- button -- ground
                    \end{verbatim}
                \item \textbf{Pull down resistor}:  A pull-down resistor connects the input pin directly to ground (0V) It holds the pin at a logical LOW (0) when a switch is open
                    \bigbreak \noindent 
                    When you close the switch, it connects the pin to the voltage supply, overriding the resistor and pulling the pin to a logical HIGH (
            \end{itemize}
            Many MCUs contain configurable internal pull-up and pull-down resistors, so an external resistor is not always necessary.
        \item \textbf{The pin is not the same as the GPIO peripheral}: A physical MCU pin is connected to internal routing circuitry. That circuitry determines which internal component controls the pin.
            \bigbreak \noindent 
            A pin might be routed to:
            \begin{itemize}
                \item GPIO
                \item UART transmit or receive
                \item SPI clock or data
                \item I$^2$C
                \item A timer PWM channel
                \item An ADC
                \item Another peripheral
            \end{itemize}
            Therefore, a pin marked PA5 does not necessarily have to act as an ordinary GPIO. It could be assigned an alternate function
        \item \textbf{Alternate function}: An alternate function on a microcontroller (MCU) is a special feature that lets a general-purpose input/output (GPIO) pin connect to an internal hardware peripheral instead of just acting as a basic digital input or output
        \item \textbf{STM32 GPIO naming}: 
            \begin{itemize}
                \item PA5 means port A, pin 5.
                \item PB3 means port B, pin 3.
                \item PC13 means port C, pin 13.
            \end{itemize}
            GPIO pins are organized into ports, usually containing up to 16 pins:
            \begin{itemize}
                \item \textbf{GPIOA:} PA0 through PA15
                \item \textbf{GPIOB:} PB0 through PB15
                \item \textbf{GPIOC:} PC0 through PC15
            \end{itemize}
        \item \textbf{GPIO modes}: An STM32 GPIO pin commonly has four principal modes.
            \begin{center}
                \begin{tabularx}{\textwidth}{@{}XX@{}}
                    \toprule
                    \textbf{Mode}&	\textbf{Purpose} \\
                    \midrule
                    Input	&Read an external digital signal\\[2ex]
                    Output	&Drive the pin high or low\\[2ex]
                    Alternate function	&Connect the pin to UART, SPI, timers, etc.\\[2ex]
                    Analog	&Connect the pin to analog circuitry or reduce digital power consumption\\
                    \bottomrule
                \end{tabularx}
            \end{center}
            Selecting the correct mode is essential. Writing an output value does not make a pin behave as an output if it is still configured as an input.
        \item \textbf{Output driver types}: When operating as an output, a GPIO pin can commonly use either push-pull or open-drain behavior.
            \bigbreak \noindent 
            A push-pull output can actively drive both voltage levels:
            \bigbreak \noindent 
            \begin{itemize}
                \item Write 1 $\to$ pin driven high
                \item Write 0 $\to$ pin driven low
            \end{itemize}
            This is normally used for:
            \begin{itemize}
                \item LEDs
                \item Digital control signals
                \item SPI signals
                \item General output pins
            \end{itemize}
        \item \textbf{Open-drain}: An open-drain output is a digital pin configuration that can pull a signal down to ground or disconnect completely, but it cannot actively drive a high voltage on its own
            \bigbreak \noindent 
            If any device pulls the line low, the line is low. The line becomes high only when every device releases it.
        \item \textbf{Output speed}: STM32 GPIO configuration includes an output speed setting. This does not normally mean how frequently software toggles the pin. It primarily controls how quickly the output voltage transitions between low and high.
            \bigbreak \noindent 
            A faster edge supports higher-frequency signals but can cause:
            \begin{itemize}
                \item More electromagnetic interference
                \item More power consumption
                \item More ringing and signal-integrity problems
            \end{itemize}
            For a slowly blinking LED, a low-speed setting is sufficient.
        \item \textbf{GPIO registers}: GPIO is controlled through memory-mapped registers. These registers occupy addresses in the MCU’s memory map, but accessing them controls hardware rather than ordinary RAM.
            \bigbreak \noindent 
            For an STM32 GPIO port, important registers commonly include:
            \begin{center}
                \begin{tabularx}{\textwidth}{@{}XX@{}}
                    \toprule
                    \textbf{Register}&	\textbf{Purpose} \\
                    \midrule
                    MODER	&Select input, output, alternate-function, or analog mode\\[2ex]
                    OTYPER	&Select push-pull or open-drain output\\[2ex]
                    OSPEEDR	&Select output transition speed\\[2ex]
                    PUPDR	&Configure pull-up or pull-down resistors\\[2ex]
                    IDR	&Read current input states\\[2ex]
                    ODR	&Read or change output states\\[2ex]
                    BSRR	&Atomically set or reset output bits\\[2ex]
                    AFRL, AFRH	&Select alternate peripheral functions \\
                    \bottomrule
                \end{tabularx}
            \end{center}
            Each pin is represented by one or more bits in these registers.
        \item \textbf{Configuring a pin as an output}: Suppose we want to configure PA5 as an output.
            \bigbreak \noindent 
            On STM32 devices, each pin has a two-bit field in MODER:
            \begin{center}
                \begin{tabularx}{\textwidth}{@{}XX@{}}
                    \toprule
                    \textbf{MODER field} &\textbf{Mode} \\
                    \midrule
                    00	&Input\\[2ex]
                    01	&Output\\[2ex]
                    10	&Alternate function\\[2ex]
                    11	&Analog \\
                    \bottomrule
                \end{tabularx}
            \end{center}
            Because PA5 is pin 5, its mode is controlled by bits 11 ($2(5)+1$) and 10 ($2(5)$):
            \bigbreak \noindent 
            A simplified register-level configuration looks like:
            \bigbreak \noindent 
            \begin{cppcode}
                GPIOA->MODER &= ~(3U << (5U * 2U));  // Clear both mode bits
                GPIOA->MODER |=  (1U << (5U * 2U));  // Select output mode
            \end{cppcode}
            The first operation clears the two-bit field. The second sets it to 01.
            \bigbreak \noindent 
            You could modify the output data register:
            \bigbreak \noindent 
            \begin{cppcode}
                GPIOA->ODR |= (1U << 5);    // Set PA5 high
                GPIOA->ODR &= ~(1U << 5);   // Set PA5 low
            \end{cppcode}
            \bigbreak \noindent 
            However, STM32 generally provides the BSRR register for safer individual-bit changes:
            \bigbreak \noindent 
            \begin{cppcode}
                GPIOA->BSRR = (1U << 5);         // Set PA5 high
                GPIOA->BSRR = (1U << (5 + 16));  // Set PA5 low
            \end{cppcode}
            In BSRR:
            \begin{itemize}
                \item Bits 0–15 set corresponding pins.
                \item Bits 16–31 reset corresponding pins.
            \end{itemize}
            This avoids a read-modify-write sequence and makes the operation atomic with respect to interrupts.
            \bigbreak \noindent 
            BSRR means Bit Set/Reset Register. It is a 32-bit, write-only register divided into two halves. The important feature is that each operation requires only one write to the peripheral.
            \bigbreak \noindent 
            After the CPU writes to GPIOA->BSRR, the write travels through the STM32 bus system to the GPIOA peripheral. The BSRR does not store the value like an ordinary variable; GPIOA interprets it as a set/reset command.
        \item \textbf{Reading an input}: To test whether PC13 is high:
            \bigbreak \noindent 
            \begin{cppcode}
                if (GPIOC->IDR & (1U << 13)) {
                    // PC13 is high
                } else {
                    // PC13 is low
                }
            \end{cppcode}
            IDR contains one bit for each pin in the port. Reading IDR observes the pin’s actual digital input state. Reading ODR observes the value stored in the output latch. These values can differ under certain electrical conditions.
        \item \textbf{RCC}: RCC stands for Reset and Clock Control, a dedicated hardware register block inside STM32 microcontrollers that manages system resets and distributes clock signals
        \item \textbf{STM32 bus architecture}:  STM32 microcontrollers control internal system buses and attached hardware via memory-mapped control, clock, and configuration registers belonging to the Reset and Clock Control (RCC) subsystem and individual peripheral blocks
            \bigbreak \noindent 
            STM32 chips (such as the STM32F4/F7 series) use an AHB/APB bus matrix to connect the CPU and DMA masters to memory and peripherals
            \begin{itemize}
                \item \textbf{AHB (Advanced High-performance Bus):} Used for high-speed components like Flash, SRAM, and DMA. Divided into AHB1 and AHB2 for core/GPIO acces
                \item \textbf{APB (Advanced Peripheral Bus):} Used for lower-speed peripherals. Divided into APB1 (timers, I2C, SPI2/3, UART2-5) and APB2 (high-speed SPI1, USART1, ADCs)
            \end{itemize}
            Peripherals on these buses are disabled by default after a reset to save power. You must enable their respective bus clocks using RCC registers before reading or writing to any peripheral register
            \bigbreak \noindent 
            \textbf{Note:} ENR stands for enable register
        \item \textbf{Enabling the GPIO peripheral clock}: On many MCUs, a GPIO peripheral does nothing until its clock is enabled. This reduces power consumption for unused hardware.
            \bigbreak \noindent 
            On the STM32F446RE, GPIO ports are connected to the AHB1 (advanced high performance bus 1) peripheral bus. GPIOA can be enabled with:
            \bigbreak \noindent 
            \begin{cppcode}
            RCC->AHB1ENR |= RCC_AHB1ENR_GPIOAEN;
            \end{cppcode}
            The sequence for a simple output is therefore:
            \begin{enumerate}
                \item Enable the GPIO port’s clock.
                \item Configure the pin mode.
                \item Configure its output type, speed, and pull resistors.
                \item Write the desired output state.
            \end{enumerate}
            \bigbreak \noindent 
            \begin{cppcode}
                // Enable GPIO port A clock
                RCC->AHB1ENR |= RCC_AHB1ENR_GPIOAEN;

                // Clear port A pin 5
                GPIOA->MODER &= ~(3U << (5U * 2U));

                // Set port A pin 5 output
                GPIOA->MODER |=  (1U << (5U * 2U));

                GPIOA->OTYPER &= ~(1U << 5);  // Push-pull
                GPIOA->PUPDR  &= ~(3U << (5U * 2U));

                // Set high
                GPIOA->BSRR = (1U << 5);      // PA5 high
            \end{cppcode}
        \item \textbf{Electrical limitations}: A GPIO pin is a small digital interface, not a general-purpose power source.
            \bigbreak \noindent 
            Important restrictions include:
            \begin{itemize}
                \item Maximum voltage on the pin
                \item Maximum current per pin
                \item Maximum total current for a port or the MCU
                \item Whether the pin is 5 V tolerant
                \item Whether analog functionality is supported
                \item Whether the pin is reserved for debugging or another board function
            \end{itemize}
        \item \textbf{GPIO and interrupts}: Instead of repeatedly checking an input:
            \bigbreak \noindent 
            \begin{cppcode}
                while (1) {
                    if (GPIOC->IDR & (1U << 13)) {
                        // Handle input
                    }
                }
            \end{cppcode}
            \bigbreak \noindent 
            the MCU can configure the pin to generate an interrupt when the signal changes.
            \bigbreak \noindent 
            Common trigger choices include:
            \begin{itemize}
                \item Rising edge: low to high
                \item Falling edge: high to low
                \item Both edges
            \end{itemize}
        \item \textbf{Voltage labels}: 
            \begin{itemize}
                \item \(V_{DD}\) represents the positive supply voltage for a circuit, particularly those using field-effect transistors (FETs) like CMOS chips. The abbreviation stands for Voltage at the Drain
                \item \(V_{SS}\) (Voltage at the Source): The negative supply or ground rail paired with \(V_{DD}\) in FET/CMOS circuits.
                \item \(V_{CC}\) / \(V_{EE}\): The equivalents used in bipolar junction transistor (BJT) circuits. \(V_{CC}\) stands for Voltage at the Collector (positive), and \(V_{EE}\) stands for Voltage at the Emitter (negative/ground).
            \end{itemize}
        \item \textbf{More on push-pull}: A GPIO pin configured as an output must electrically drive the external wire. The two common output-driver designs are push-pull and open-drain.


    \end{itemize}

    \pagebreak 
    \subsection{MCU Architecture}
    \begin{itemize}
        \item \textbf{Intro}: A Microcontroller Unit (MCU) architecture integrates a central processing unit, memory, and input/output peripherals onto a single integrated circuit
            \begin{itemize}
                \item \textbf{CPU (Central Processing Unit):} The brain that fetches, decodes, and executes program instructions. It includes the Arithmetic Logic Unit (ALU) for math and logic, and the Control Unit for sequencing operations
                \item \textbf{Program memory (ROM / flash)}: Non-volatile storage that holds your permanent application code so it is not lost when power turns off. 
                \item \textbf{Data Memory (RAM):} Volatile storage used by the CPU to temporarily save active variables and runtime data
                \item \textbf{System Bus:} An internal highway of wires that lets the CPU talk to memory and hardware peripherals
                \item \textbf{Peripherals:} Built-in hardware tools that let the chip interact with the physical world. Common examples include:
                    \begin{itemize}
                        \item \textbf{GPIO:} General-purpose input/output pins for turning things on or off.
                        \item \textbf{ADC:} Analog-to-digital converters to read voltages from sensors.
                        \item \textbf{Timers/Counters:} For measuring time, creating delays, or generating Pulse-Width Modulation (PWM) signals.
                        \item \textbf{Communication Interfaces:} Protocols like UART, SPI, and I$^2$C to talk to other chips
                    \end{itemize}
            \end{itemize}
        \item \textbf{Von Neumann architecture}: Uses a single shared memory space and bus for both data and instructions. It is simple to design but creates a data bottleneck because the CPU cannot read an instruction and read/write data at the exact same time
        \item \textbf{Harvard architecture}: Uses physically separate memory blocks and separate buses for code and data. This lets the processor fetch an instruction while simultaneously accessing data, which speeds up real-time execution. Most modern MCUs use variations of this model
        \item \textbf{Execute-in-place (XiP)}:  instructions do not necessarily need to move from ROM to RAM before execution. Most modern microcontrollers execute code directly from flash memory (ROM) using a technique called Execute-in-Place (XiP).
            \bigbreak \noindent 
            In most microcontrollers (like ARM Cortex-M, AVR, or PIC), the CPU is wired to fetch instructions directly from the Flash/ROM bus.
            \bigbreak \noindent 
            Microcontrollers have very limited RAM (often only kilobytes) but larger Flash memory. Keeping the code in Flash saves precious RAM for data and variables.
            \bigbreak \noindent 
            Flash memory is generally slower than RAM. For ultra-high-speed operations, Flash access time can become a performance bottleneck.
            \bigbreak \noindent 
            If instructions are fetched from Flash while data is simultaneously read or written in RAM, that is a classic implementation of the Harvard architecture.
        \item \textbf{Flash memory}: Flash holds data without power. It normally contains:
            \begin{itemize}
                \item The vector table
                \item Machine instructions
                \item Read-only constants
                \item Initial values for initialized global variables
            \end{itemize}
            On the STM32F446, main flash begins at:
            \bigbreak \noindent 
            \begin{ccode}
            0x08 00 00 00
            \end{ccode}
            \bigbreak \noindent 
            The processor normally executes instructions directly from flash. This is commonly called execute in place.
            \bigbreak \noindent 
            Flash is persistent but slower and more complicated to modify than SRAM. You cannot generally treat it as ordinary writable memory.
        \item \textbf{SRAM}: SRAM is volatile working memory. It loses its contents when power is removed.
            \bigbreak \noindent 
            It contains things such as:
            \begin{itemize}
                \item Global and static variables
                \item The stack
                \item The heap, if one is used
                \item Communication buffers
                \item Temporary computation results
            \end{itemize}
            The primary SRAM address range begins at:
            \bigbreak \noindent 
            \begin{ccode}
            0x20000000
            \end{ccode}
            \bigbreak \noindent 
            \begin{center}
                \begin{tabularx}{\textwidth}{@{}XX@{}}
                    \toprule
                    \textbf{Address range}&	\textbf{Purpose} \\
                    \midrule
                    0x00000000...	&Boot-memory alias\\[2ex]
                    0x08000000–0x0807FFFF	&512 KiB internal Flash\\[2ex]
                    Mostly reserved/unmapped space	&No physical memory responds here\\[2ex]
                    Near 0x1FFF0000	&Factory system memory and bootloader\\[2ex]
                    0x20000000–0x2001FFFF	&128 KiB SRAM \\
                    \bottomrule
                \end{tabularx}
            \end{center}
            The address space is not one enormous physical memory device. Large address ranges are simply unused or reserved.
        \item \textbf{The boot-memory alias}: 0x00000000 is particularly important. At reset, the Cortex-M4 expects the vector table to appear there:
            \begin{center}
                \begin{tabularx}{\textwidth}{@{}XX@{}}
                    \toprule
                    \textbf{Address}&	\textbf{Initial meaning} \\
                    \midrule
                    0x00000000	&Initial Main Stack Pointer value\\[2ex]
                    0x00000004	&Address of the reset handler\\[2ex]
                    0x00000008	&Address of the NMI handler\\[2ex]
                    0x0000000C	&Address of the HardFault handler \\
                    \bottomrule
                \end{tabularx}
            \end{center}
        \item \textbf{Programs}: In ordinary embedded development, Flash contains one firmware image:
            \begin{verbatim}
                Flash
                |-- Vector table
                |-- Program instructions
                |-- Read-only constants
                |-- Initial values for global variables
            \end{verbatim}
            \bigbreak \noindent 
            That firmware might contain many modules, functions, tasks, or an RTOS, but collectively it is normally considered one program.
            \bigbreak \noindent 
            You can store multiple independent firmware images if you create a bootloader and partition Flash.
            \begin{center}
                \begin{tabularx}{\textwidth}{@{}XX@{}}
                    \toprule
                    \textbf{Region}&	\textbf{Example size} \\
                    \midrule
                    Bootloader	&32 KiB\\[2ex]
                    Application A	&200 KiB\\[2ex]
                    Application B	&200 KiB\\[2ex]
                    Configuration/metadata	&Remaining space \\
                    \bottomrule
                \end{tabularx}
            \end{center}
            The bootloader would select an application, set the application’s vector-table location, load its initial stack pointer, and jump to its reset handler.
            \bigbreak \noindent 
            Without a custom bootloader: normally one firmware program
        \item \textbf{Address space}: On the STM32F446RE, the address space is primarily a physical address space, not a virtual address space.
            \bigbreak \noindent 
            When the Cortex-M4 executes:
            \bigbreak \noindent 
            \begin{nasmcode}
                LDR r0, [r1]
            \end{nasmcode}
            \bigbreak \noindent 
            and r1 contains 0x20000100, that address is placed onto the MCU’s internal bus. The bus interconnect routes it directly to the SRAM hardware.
            \bigbreak \noindent 
            Memory on the MCU is usually not virtual
        \item \textbf{Startup chain}: 
            \begin{enumerate}
                \item Reset happens 
                \item boot mode selects memory 
                \item CPU reads the vector table 
                \item CPU enters the reset handler 
                \item startup code calls main()
            \end{enumerate}
        \item \textbf{Resets}: A reset returns the processor and much of the microcontroller hardware to a known initial state. It does not erase the program stored in Flash.
            \bigbreak \noindent 
            A reset may be caused by:
            \begin{itemize}
                \item Power being applied
                \item The NRST reset pin being asserted
                \item A watchdog timeout
                \item Brownout or insufficient supply voltage
                \item Software requesting a reset
                \item A debugger pressing “reset”
            \end{itemize}
            Different reset sources affect slightly different parts of the MCU, but they generally restart program execution using the same boot process. 
            \bigbreak \noindent 
            After reset:
            \begin{itemize}
                \item The CPU registers return to defined reset states.
                \item Most peripherals return to their reset states.
                \item The CPU uses the Main Stack Pointer, or MSP.
                \item The CPU begins in privileged Thread mode.
                \item Flash retains its contents.
                \item Ordinary SRAM contents should be considered undefined, even if some electrical values happen to remain.
            \end{itemize}
        \item \textbf{Booting}: Booting is the process of deciding which code the processor should begin executing and then preparing that code to run.
            \bigbreak \noindent 
            For a normal embedded application, the process is:
            \begin{enumerate}
                \item Reset
                \item Select boot memory
                \item Read initial stack pointer
                \item Read reset-handler address
                \item Execute Reset\_Handler
                \item Initialize RAM and hardware
                \item Call main()
            \end{enumerate}
            The STM32F446 can boot from different physical memory regions. Boot configuration pins and option settings determine which region appears at address 0x00000000.
            \bigbreak \noindent 
            The primary choices are:
            \begin{center}
                \begin{tabularx}{\textwidth}{@{}XX@{}}
                    \toprule
                    \textbf{Boot source}&	\textbf{Purpose} \\
                    \midrule
                    Main Flash	& Run your normal firmware\\[2ex]
                    System memory&	Run ST’s factory-programmed bootloader\\[2ex]
                    SRAM	&Run a program previously loaded into SRAM \\
                    \bottomrule
                \end{tabularx}
            \end{center}
            The Nucleo board normally holds BOOT0 low, selecting the main Flash. 
            \bigbreak \noindent 
            Your program physically begins at:
            \bigbreak \noindent 
            \begin{bashcode}
            0x08000000
            \end{bashcode}
            \bigbreak \noindent 
            During normal Flash boot, the STM32 aliases this memory at:
            \bigbreak \noindent 
            \begin{bashcode}
            0x00000000
            \end{bashcode}
            \bigbreak \noindent 
            Both addresses refer to the same beginning of Flash during this boot configuration. The CPU reads the initial startup information from 0x00000000, while your linker normally places it physically at 0x08000000.
        \item \textbf{System bootloader}: STM32 contains factory-programmed code called the system bootloader. It can be used to program Flash through supported communication interfaces without using the normal debugger.
            \bigbreak \noindent 
            This factory bootloader is separate from your application and is not erased when you program ordinary Flash.
            \bigbreak \noindent 
            An MCU system bootloader is a dedicated piece of firmware stored in a device's non-volatile memory that executes immediately upon power-up or reset to initialize hardware and manage application loading or updates
            \bigbreak \noindent 
            Configures essential low-level hardware like clocks, memory regions, and input/output ports before handing control over
            \bigbreak \noindent 
            Usually sits at a fixed address range completely distinct from the main user flash memory (for example, starting at 0x1FFF0000 on certain STM32 chips).
            \bigbreak \noindent 
            This memory is write-protected at the factory level and cannot be erased or modified by your software. It is always available as a fallback recovery mechanism.
            \bigbreak \noindent 
            You activate the ST system bootloader when you want to program or recover the MCU without running the firmware currently stored in main flash.
        \item \textbf{SRAM boot}: The CPU can also start from SRAM. This is primarily useful for debugging, testing, or specialized loaders.
            \bigbreak \noindent 
            Because SRAM is volatile, the program must already have been loaded into SRAM by something such as a debugger.
        \item \textbf{Vector table}:  The vector table is an array of 32-bit values that tells the Cortex-M4:
            \begin{enumerate}
                \item What the initial stack pointer should be.
                \item Where the reset handler is.
                \item Where every exception and interrupt handler is.
            \end{enumerate}
            It is called a “vector” table because each entry directs, or vectors, execution to a handler. For a normal application, the vector table is placed at the beginning of Flash:
            \bigbreak \noindent 
            The vector table and your program at 0x0800.. do not clash because the vector table is part of your program’s flash image. When we say the program “starts at 0x08000000,” that does not mean the first machine instruction is stored there.
            \begin{center}
                \begin{tabularx}{\textwidth}{@{}XX@{}}
                    \toprule
                    \textbf{Address}&	\textbf{Contents} \\
                    \midrule
                    0x08000000	&Initial Main Stack Pointer value\\[2ex]
                    0x08000004	&Address of Reset\_Handler\\[2ex]
                    0x08000008	&Address of NMI\_Handler\\[2ex]
                    0x0800000C	&Address of HardFault\_Handler\\[2ex]
                    ...	&Remaining exception and interrupt vectors\\[2ex]
                    After vector table	&Startup code, functions, constants, and other program data \\[2ex]
                    Later address	&main() machine code \\
                    \bottomrule
                \end{tabularx}
            \end{center}
            Each independently bootable firmware program normally has its own vector table. The vector table belongs to the firmware image and tells the processor:
            \begin{itemize}
                \item Where its stack begins
                \item Where its reset code begins
                \item Where its exception and interrupt handlers are
            \end{itemize}
        \item \textbf{Multiple programs in flash}: Suppose you have a bootloader and an application:
            \bigbreak \noindent 
            \begin{ccode}
            0x08000000  Bootloader vector table
                Bootloader code

            0x08020000  Application vector table
                Application code
            \end{ccode}
            \bigbreak \noindent 
            At reset, the processor uses the bootloader’s vector table because it is at the boot address. When the bootloader wants to run the application, it reads the application’s first two vector entries:
            \bigbreak \noindent 
            \begin{ccode}
                uint32_t app_msp   = *(uint32_t *)0x08020000;
                uint32_t app_reset = *(uint32_t *)0x08020004;
            \end{ccode}
            \bigbreak \noindent 
            It then approximately performs:
            \bigbreak \noindent 
            \begin{ccode}
                SCB->VTOR = 0x08020000;  // Use application's vector table
                __set_MSP(app_msp);      // Use application's stack
                jump_to(app_reset);      // Run application's Reset_Handler
            \end{ccode}
            \bigbreak \noindent 
            Therefore, multiple firmware images can coexist, but only one vector table is active for interrupt dispatch at a time.
        \item \textbf{The Main Stack Pointer}: The Main Stack Pointer, or MSP, is a CPU register containing the current top of the main stack.
            \bigbreak \noindent 
            On Cortex-M processors, register R13 serves as the stack pointer. There are actually two possible stack pointers:
            \begin{itemize}
                \item \textbf{MSP:} Main Stack Pointer
                \item \textbf{PSP:} Process Stack Pointer
            \end{itemize}
            The MSP is normally used:
            \begin{itemize}
                \item Immediately after reset
                \item During startup
                \item By exception and interrupt handlers
                \item Often by simple bare-metal applications for everything
            \end{itemize}
            The PSP is optional and is commonly used by an operating system to give application threads separate stacks.
            \bigbreak \noindent 
            The first vector-table entry contains the initial value that the processor loads into the MSP:
            \bigbreak \noindent 
            \begin{ccode}
                Address       Value          Meaning
                0x08000000    0x20020000     Initial MSP
                0x08000004    0x08000151     Reset_Handler address
            \end{ccode}
            \bigbreak \noindent 
            During reset, the processor effectively performs:
            \bigbreak \noindent 
            \begin{ccode}
                MSP <- memory[0x00000000]
                PC  <- memory[0x00000004]
            \end{ccode}
            \bigbreak \noindent 
            When booting from main flash, those locations correspond to:
            \bigbreak \noindent 
            \begin{ccode}
                MSP <- memory[0x08000000]
                PC  <- memory[0x08000004]
            \end{ccode}
            \begin{itemize}
                \item 0x08000000 is where the initial MSP value is stored in flash.
                \item 0x20020000 is the SRAM address loaded into the MSP register.
            \end{itemize}
            \textbf{Note:} Stacks on Cortex-M normally grow toward lower addresses.
        \item \textbf{Peripherals}: A peripheral is a hardware unit that performs a specialized task.
            \begin{itemize}
                \item \textbf{GPIO:} reads and controls pins
                \item \textbf{USART/UART:} serial communication
                \item \textbf{Timers:} timing, counting, PWM and input capture
                \item \textbf{ADC:} converts analog voltage into a number
                \item \textbf{DMA:} moves data without making the CPU copy every element
                \item \textbf{RCC:} controls clocks and resets
                \item \textbf{EXTI:} detects external pin events
            \end{itemize}
            A peripheral normally contains configuration registers, status registers and data registers.
        \item \textbf{Memory-mapped I/O}: The Cortex‑M4 communicates with peripherals using ordinary memory addresses.
            \bigbreak \noindent 
            \begin{ccode}
            uint32_t value = *(volatile uint32_t *)0x40020010;
            \end{ccode}
            \bigbreak \noindent 
            The C expression appears to read memory, but address 0x40020010 belongs to a GPIO peripheral. The access therefore reads a hardware register rather than SRAM.
            \bigbreak \noindent 
            When the CPU tries to access an address assigned to a peripheral, the hardware bus or memory management unit routes the request directly to the device's physical registers rather than to system RAM
            \bigbreak \noindent 
            A peripheral register can change for reasons outside the current C code. The compiler must perform each requested access rather than caching or eliminating it. Therefore register pointers are normally declared volatile:
            \bigbreak \noindent 
            \begin{ccode}
                volatile uint32_t *status =
                    (volatile uint32_t *)0x40020010;
            \end{ccode}
        \item \textbf{Registers and special-function registers}: The word “register” has two related meanings in embedded programming.
            \begin{itemize}
                \item \textbf{CPU registers}: These are storage locations directly inside the Cortex‑M4 core:
                    \begin{itemize}
                        \item R0, R1, R2, ... R15
                        \item xPSR
                        \item CONTROL
                        \item PRIMASK
                        \item BASEPRI
                        \item FAULTMASK
                    \end{itemize}
                    CPU registers are operated on directly by machine instructions.
                \item \textbf{Peripheral registers}: These are memory-mapped hardware control locations:
                    \begin{itemize}
                        \item GPIOA\_MODER
                        \item GPIOA\_IDR
                        \item GPIOA\_ODR
                        \item RCC\_AHB1ENR
                        \item TIM2\_CR1
                        \item USART2\_DR
                    \end{itemize}
                    They are sometimes called special-function registers, or SFRs.
            \end{itemize}
        \item \textbf{Common register types}
            \begin{center}
                \begin{tabularx}{\textwidth}{@{}XX@{}}
                    \toprule
                    \textbf{Type}&	\textbf{Behavior} \\
                    \midrule
                    Read/write	&Software can read and modify it\\[2ex]
                    Read-only	&Reports hardware status\\[2ex]
                    Write-only	&Writing commands the hardware\\[2ex]
                    Write-one-to-clear	&Writing 1 clears a status flag\\[2ex]
                    Set/reset register	&Different bits atomically set or clear outputs \\
                    \bottomrule
                \end{tabularx}
            \end{center}
            This is why blindly applying code such as reg |= mask can be incorrect. The register’s documented semantics must be checked first.
        \item \textbf{Cortex‑M4 register set}
            \begin{center}
                \begin{tabularx}{\textwidth}{@{}XX@{}}
                    \toprule
                    \textbf{Register}&	\textbf{Conventional role} \\
                    \midrule
                    R0–R3	&Arguments, return values and temporary values\\[2ex]
                    R4–R11	&General-purpose, normally preserved by called functions\\[2ex]
                    R12	&Temporary scratch register\\[2ex]
                    R13	&Stack pointer\\[2ex]
                    R14	&Link register\\[2ex]
                    R15	&Program counter \\
                    \bottomrule
                \end{tabularx}
            \end{center}
            R14, is a dedicated core register in the ARM Cortex-M4 processor that holds the return address for function calls, subroutines, and exception handlers. 
            \bigbreak \noindent 
            The calling convention, rather than the hardware alone, gives many of these registers their software roles.
        \item \textbf{Program status registers}: Several status components are collectively accessed as xPSR:
            \begin{itemize}
                \item \textbf{APSR:} arithmetic flags
                \item \textbf{IPSR:} current exception number
                \item \textbf{EPSR:} execution-state information
            \end{itemize}
            \begin{center}
                \begin{tabularx}{\textwidth}{@{}XX@{}}
                    \toprule
                    \textbf{Flag}&	\textbf{Meaning} \\
                    \midrule
                    N	&Negative result\\[2ex]
                    Z	&Zero result\\[2ex]
                    C	&Carry or unsigned condition\\[2ex]
                    V	&Signed overflow\\[2ex]
                    Q	&Saturation occurred \\
                    \bottomrule
                \end{tabularx}
            \end{center}
        \item \textbf{Reset sequence}: Reset is the bridge between powered hardware and your C main() function.
            \bigbreak \noindent 
            \fig{.2}{./figures/1.png}
            \bigbreak \noindent 
            More precisely:
            \begin{enumerate}
                \item The processor enters reset state.
                \item The vector table is made visible at the boot address.
                \item The processor reads vector entry 0 and loads MSP.
                \item It reads vector entry 1 and loads the reset-handler address into PC.
                \item Execution begins in Reset\_Handler.
                \item Startup code copies initialized data from flash into SRAM.
                \item Startup code clears the .bss section.
                \item System initialization may configure clocks and the FPU.
                \item C++ constructors are run when applicable.
                \item main() is called.
            \end{enumerate}
        \item \textbf{Why .data must be copied}: Consider
            \bigbreak \noindent 
            \begin{ccode}
            int counter = 5;
            \end{ccode}
            \bigbreak \noindent 
            counter must be writable, so its runtime location is in SRAM. But SRAM loses its contents when powered off. Therefore the initial value 5 is stored in flash and copied into the SRAM location during startup.
        \item \textbf{Why .bss must be cleared}: Consider
            \bigbreak \noindent 
            \begin{ccode}
            static int errors;
            \end{ccode}
            \bigbreak \noindent 
            C requires errors to begin as zero. Instead of storing thousands of zero bytes in flash, the linker places it in .bss, and startup code clears that memory range.
        \item \textbf{Privileged and unprivileged execution}: Cortex‑M4 execution has two modes:
            \begin{itemize}
                \item \textbf{Thread mode:} ordinary program or task execution
                \item \textbf{Handler mode:} exception and interrupt handlers
            \end{itemize}
            Thread mode may be privileged or unprivileged. Handler mode is always privileged
            \bigbreak \noindent 
            Privileged software can perform operations such as:
            \begin{itemize}
                \item Configure the MPU
                \item Change interrupt masks
                \item Access restricted memory regions
                \item Modify important system-control registers
            \end{itemize}
            Unprivileged software is restricted.
            \bigbreak \noindent 
            The CONTROL.nPRIV bit controls Thread-mode privilege:
            \bigbreak \noindent 
            \begin{ccode}
                0 = privileged
                1 = unprivileged
            \end{ccode}
            \bigbreak \noindent 
            Privileged Thread mode can drop into unprivileged mode directly. Unprivileged Thread mode generally cannot restore its own privilege. It must request a privileged exception handler, commonly through SVC.
        \item \textbf{Exceptions and interrupts}: An exception is any event that causes the processor to suspend normal instruction execution and run a handler.
            \bigbreak \noindent 
            An interrupt is generally an asynchronous, hardware-generated kind of exception
        \item \textbf{Synchronous exceptions}: These result from the instruction being executed:
            \begin{itemize}
                \item Undefined instruction
                \item Invalid memory access
                \item Divide-by-zero trap, when enabled
                \item Supervisor call through SVC
            \end{itemize}
        \item \textbf{Asynchronous interrupts}: These arrive from hardware independently of the current instruction:
            \begin{itemize}
                \item Timer expiration
                \item UART byte received
                \item GPIO input edge
                \item DMA transfer completed
                \item ADC conversion completed
            \end{itemize}
        \item \textbf{Exception entry}: When an exception is accepted, the processor automatically saves a basic context frame containing:
            \bigbreak \noindent 
            \begin{ccode}
                R0
                R1
                R2
                R3
                R12
                LR
                PC
                xPSR
            \end{ccode}
            \bigbreak \noindent 
            It then:
            \begin{enumerate}
                \item Selects MSP for Handler mode.
                \item Obtains the handler address from the vector table.
                \item Updates the current exception state.
                \item Begins executing the handler.
            \end{enumerate}
            On exception return, the hardware restores the saved registers and resumes the interrupted code.
        \item \textbf{Priorities and nesting}: A higher-priority interrupt can preempt a lower-priority handler. On Cortex‑M, a lower numerical priority value represents a higher logical priority.
            \bigbreak \noindent 
            \begin{ccode}
                priority 0  = highest
                priority 3  = higher than priority 8
                priority 15 = lower
            \end{ccode}
            \bigbreak \noindent 
            The STM32F446 implements a finite number of priority bits, so not every bit in an 8-bit priority field is meaningful.
        \item \textbf{CMSIS}: CMSIS (Cortex Microcontroller Software Interface Standard) is a vendor-independent hardware abstraction layer developed by Arm that provides consistent software interfaces for Arm Cortex-M processor-based microcontrollers like the STM32.
            \bigbreak \noindent 
            CMSIS assigns negative interrupt numbers to core exceptions and nonnegative numbers to device-specific interrupts
        \item \textbf{Interrupt Service Routine (ISR)}: Another name for interrupt handler 
        \item \textbf{NVIC interrupt controller}: NVIC means Nested Vectored Interrupt Controller. It is part of the Cortex‑M core-peripheral architecture, not an ordinary STM32 timer or GPIO peripheral.
            \bigbreak \noindent 
            The NVIC (Nested Vectored Interrupt Controller) is a dedicated hardware block inside ARM Cortex-M based STM32 microcontrollers that manages, prioritizes, and processes all asynchronous system exceptions and hardware interrupts
            \begin{itemize}
                \item \textbf{Interrupt Prioritization:} Decides which task gets attention first; lower numerical values mean higher urgency (Priority 0 is the highest).
                \item \textbf{Preemption and Nesting:} Allows higher-priority interrupts to interrupt (nest inside) lower-priority active interrupt service routines (ISRs). 
                \item \textbf{Vector Table Lookup:} Directly points the CPU to the exact memory address of the correct Interrupt Service Routine (ISR) when a signal triggers. 
                \item \textbf{Tail-Chaining:} Optimizes back-to-back interrupts by skipping stack pop/push operations to reduce latency.
            \end{itemize}
            \bigbreak \noindent 
            The NVIC records whether each external interrupt is:
            \begin{itemize}
                \item Enabled
                \item Pending
                \item Active
                \item Assigned a particular priority
            \end{itemize}
        \item \textbf{TIM2}: STM32 TIM2 is a heavy-duty, general-purpose 32-bit timer found on many STM32 microcontrollers
        \item \textbf{Peripheral versus NVIC responsibilities}: Suppose TIM2 reaches its programmed count:
            \begin{enumerate}
                \item TIM2 sets its own update flag.
                \item If the timer’s update interrupt is enabled, it signals the NVIC.
                \item The NVIC marks TIM2’s interrupt pending.
                \item The CPU compares its priority with the current execution priority.
                \item The CPU enters TIM2_IRQHandler.
                \item Software clears the TIM2 update flag.
                \item The processor returns from the exception.
            \end{enumerate}
            There are therefore usually two enable controls:
            \bigbreak \noindent 
            \begin{center}
                Peripheral interrupt enable + NVIC interrupt enable
            \end{center}
            Both must be configured.
            \bigbreak \noindent 
            Similarly, clearing the NVIC pending bit alone may not solve the problem. If the peripheral’s source flag remains set, it can request the interrupt again
        \item \textbf{Interrupt enable}: An interrupt enable is a switch or control bit in a microcontroller that tells the processor whether it is allowed to respond to a specific hardware event or a global signal
        \item \textbf{SysTick timer}: 
            SysTick is a simple timer integrated with the Cortex‑M core. It is commonly used for:
            \begin{itemize}
                \item A periodic operating-system tick
                \item Millisecond timekeeping
                \item Software delays
                \item Scheduling timeouts
            \end{itemize}
            Its main registers are:
            \begin{center}
                \begin{tabularx}{\textwidth}{@{}XX@{}}
                    \toprule
                    \textbf{Register}&	\textbf{Purpose} \\
                    \midrule
                    SYST\_CSR / CTRL	&Enable, clock selection, interrupt enable and status\\[2ex]
                    SYST\_RVR / LOAD	&Reload value\\[2ex]
                    SYST\_CVR / VAL	&Current counter value\\[2ex]
                    SYST\_CALIB	&Implementation-specific calibration information \\
                    \bottomrule
                \end{tabularx}
            \end{center}
            \bigbreak \noindent 
            SysTick is convenient, but it is not a full replacement for STM32 general-purpose timers.
            \bigbreak \noindent 
            SysTick\_Config() sets the reload value, clears the current count, enables the interrupt and starts the counter.
            \bigbreak \noindent 
            Its connection to interrupts is that it can automatically generate an interrupt whenever its countdown reaches zero.
            \bigbreak \noindent 
            \begin{ccode}
                volatile uint32_t milliseconds;

                void SysTick_Handler(void)
                {
                    ++milliseconds;
                }

                int main(void)
                {
                    SysTick_Config(SystemCoreClock / 1000u);

                    while (1) {
                        // milliseconds is updated once per millisecond
                    }
                }
            \end{ccode}
            \begin{enumerate}
                \item SysTick counts downward while the processor executes main().
                \item The counter reaches zero.
                \item SysTick raises an interrupt request.
                \item The processor temporarily pauses main().
                \item The processor executes SysTick\_Handler().
                \item The handler increments milliseconds.
                \item The processor returns to the interrupted location in main()
            \end{enumerate}
        \item \textbf{Memory Protection Unit}: The MPU divides the physical address space into regions and assigns access rules.
            \bigbreak \noindent 
            A region can specify:
            \begin{itemize}
                \item Read/write permissions
                \item Privileged versus unprivileged access
                \item Executable or execute-never
                \item Memory type and caching-related attributes
                \item Region size and base address
            \end{itemize}
            Possible uses include:
            \begin{itemize}
                \item Making flash read-only
                \item Preventing instruction execution from SRAM
                \item Protecting kernel memory from application tasks
                \item Detecting stack overflows with a guard region
                \item Restricting a task’s peripheral access
            \end{itemize}
            If software violates a configured MPU rule, the processor normally raises a MemManage fault.
        \item \textbf{Floating-point unit}: The Cortex‑M4 used by the STM32F446 contains an FPv4-SP single-precision FPU.
            \bigbreak \noindent 
            Single precision means normal hardware support for 32-bit IEEE-754 float operations:
            \bigbreak \noindent 
            It does not provide the same hardware support for 64-bit double arithmetic. Double-precision operations may require software routines and can be considerably slower.
            \bigbreak \noindent 
            The floating-point register file includes registers conventionally named:
            \bigbreak \noindent 
            \begin{ccode}
            S0-S31
            \end{ccode}
            \bigbreak \noindent 
            There is also an \texttt{FPSCR} status/control register.
        \item \textbf{Floating point and interrupts}: If an interrupt handler uses the FPU, floating-point state might also need to be saved. Cortex‑M4 supports lazy floating-point context preservation, which postpones stacking that state until the handler actually executes a floating-point instruction.
            \bigbreak \noindent 
            This reduces interrupt overhead when the FPU is enabled but a particular handler does not use it.
        \item \textbf{Little-endian representation}: The STM32F446 uses little-endian byte ordering for ordinary data.
        \item \textbf{Thumb and Thumb‑2 instructions}: 




    \end{itemize}


    \pagebreak 
    \unsect{Context}
    The STMicroelectronics NUCLEO-F446RE is a highly popular, mid-range development board belonging to the STM32 Nucleo-64 ecosystem. It is designed around the powerful STM32F446RET6 microcontroller
    \bigbreak \noindent 
    Arm Cortex-M4 core processor

    \bigbreak \noindent 
    \subsection{Documents}
    \begin{itemize}
        \item \textbf{UM1724 — STM32 Nucleo-64 board user manual}:	Board connectors, jumpers, power options, LEDs, buttons, clock configuration, ST-LINK, and Arduino/Morpho pin mappings
        \item \textbf{DS10693 — STM32F446xC/E datasheet}:	Pinouts, alternate pin functions, voltage/current limits, clock limits, memory capacity, and electrical characteristics
        \item \textbf{RM0390 — STM32F446xx reference manual}:	Peripheral behavior and registers: GPIO, RCC, timers, UART, SPI, I²C, ADC, DMA, interrupts, flash, USB, and more
        \item \textbf{PM0214 — STM32 Cortex-M4 programming manual}:	CPU registers, instruction set, exception handling, NVIC, SysTick, MPU, FPU, and processor operating modes
        \item \textbf{MB1136 board schematic}:	Exact electrical connections between the STM32, LEDs, buttons, oscillators, ST-LINK, and headers
        \item \textbf{ES0298 — STM32F446 device errata}:	Known hardware limitations and their workarounds
        \item \textbf{UM1725 — STM32F4 HAL and LL driver manual}:	STM32Cube HAL and low-layer driver APIs
    \end{itemize}

    \bigbreak \noindent 
    \subsection{Packages}
    \begin{itemize}
        \item arm-none-eabi-gcc	Compiles C/C++ for the Cortex-M4
        \item arm-none-eabi-newlib	Embedded C standard library
        \item arm-none-eabi-gdb	Debugs the program running on the board
        \item openocd	Communicates with the board’s integrated ST-LINK
    \end{itemize}

    \bigbreak \noindent 
    \subsection{STM32CubeMX}
    \bigbreak \noindent 
    STM32CubeMX is a graphical hardware-configuration and code-generation program made by STMicroelectronics. It helps you configure an STM32 microcontroller without manually writing all the startup and peripheral-initialization code.

    \bigbreak \noindent 
    \subsection{Generating a Project via CLI}
    \bigbreak \noindent 
    Create a new text file named extension \texttt{.script}. Then,
    \bigbreak \noindent 
    \begin{bashcode}
        # 1. Set the target microcontroller (replace with your exact part number)
        config load STM32F466RETx
    \end{bashcode}
    \bigbreak \noindent 
    it is highly recommended to load the board configuration instead of just the MCU. Loading the board pre-configures system peripherals like the default USART pins for the built-in ST-LINK debugger, the user LEDs, and push buttons.
    \bigbreak \noindent 
    \begin{bashcode}
    board load NUCLEO-F446RE

    # 2. Setup project path and name
    project name <project name>
    project path <path>

    # 3. Choose your IDE/Toolchain (e.g., STM32CubeIDE, Makefile, EWARM, MDK-ARM)
    project toolchain Makefile

    # Saves the generated .ioc (input output configuration)
    config saveas <path>

    # 4. Generate the code and project structure
    project generate

    # 5. Exit the CLI tool
    exit
    \end{bashcode}
    \bigbreak \noindent 
    \begin{bashcode}
    stm32cubemx -q *.script
    \end{bashcode}

    \bigbreak \noindent 
    \subsection{The System}
    \begin{itemize}
        \item \textbf{Cortex-M4 core:} Executes instructions and contains the CPU registers, exception mechanism, and SysTick timer.
        \item \textbf{STM32F446RE MCU:} Combines the Cortex-M4 core with Flash, SRAM, GPIO, timers, UARTs, ADCs, and other peripherals.
        \item \textbf{NUCLEO-F446RE board:} Connects the MCU to an onboard LED, button, ST-LINK debugger, USB connector, headers, and supporting circuitry.
    \end{itemize}

    \bigbreak \noindent 
    \subsection{Basics}
    \begin{itemize}
        \item \textbf{}        


    \end{itemize}

    \bigbreak \noindent 
    \subsection{Edit, Build, Flash, Debug cycle}
    \bigbreak \noindent 
    Now that we have the project files, open the \texttt{.ioc} and verify
    \begin{itemize}
        \item PA5 is configured as GPIO\_Output.
        \item PA5 controls the green onboard LED, labeled LD2.
        \item SYS $\to$ Debug is set to Serial Wire.
        \item System Core $\to$ GPIO shows PA5 as an output.
        \item A project toolchain is selected under Project Manager.
    \end{itemize}

    In \texttt{main.c}, have
    \bigbreak \noindent 
    \begin{ccode}
        while (1)
        {
            HAL_GPIO_TogglePin(GPIOA, GPIO_PIN_5);
            HAL_Delay(500);

            /* USER CODE END WHILE */

            /* USER CODE BEGIN 3 */
        }
    \end{ccode}
    \bigbreak \noindent 
    Then, build the project with the generated makefile.
    \begin{itemize}
        \item make compiles and links.
        \item OpenOCD communicates with the onboard ST-LINK.
        \item arm-none-eabi-gdb provides source-level debugging.
    \end{itemize}
    Use \texttt{Build/*.elf} for flashing and debugging. It contains:
    \begin{itemize}
        \item Program instructions and data.
        \item Flash addresses.
        \item Function and variable names.
        \item Source-line debugging information.
    \end{itemize}
    Check the result with:
    \bigbreak \noindent 
    \begin{bashcode}
    arm-none-eabi-size build/*.elf
    \end{bashcode}
    \bigbreak \noindent 
    \begin{bashcode}
        text	  data	   bss	   dec	   hex	filename
        5544	    20	  1644	  7208	  1c28	intro.elf
    \end{bashcode}
    \bigbreak \noindent 
    Connect the USB cable to the ST-LINK USB connector and run:
    \bigbreak \noindent 
    \begin{bashcode}
    openocd -f board/st_nucleo_f4.cfg
    \end{bashcode}
    \bigbreak \noindent 
    Successful output should mention:
    \bigbreak \noindent 
    \begin{bashcode}
        STLINK
        stm32f4x.cpu
        Listening on port 3333 for gdb connections
    \end{bashcode}
    \bigbreak \noindent 
    The supplied st\_nucleo\_f4.cfg configures the ST-LINK interface, SWD transport, and STM32F4 target.
    \bigbreak \noindent 
    To flash the program, run 
    \bigbreak \noindent 
    \begin{bashcode}
    openocd \
        -f board/st_nucleo_f4.cfg \
        -c "program build/*.elf verify reset exit"
    \end{bashcode}
    \bigbreak \noindent 
    This performs four operations:
    \begin{enumerate}
        \item Erases the necessary Flash sectors.
        \item Programs the .elf
        \item Verifies the written data.
        \item Resets and starts the MCU.
    \end{enumerate}
    Your LED program should begin running immediately.
    \bigbreak \noindent 
    Debugging uses two terminals. In the first terminal, start OpenOCD:
    \bigbreak \noindent 
    \begin{bashcode}
    openocd -f board/st_nucleo_f4.cfg
    \end{bashcode}
    \bigbreak \noindent 
    Leave it running. OpenOCD exposes a GDB server on TCP port 3333. In the second terminal, start GDB:
    \bigbreak \noindent 
    \begin{bashcode}
    arm-none-eabi-gdb -tui build/intro.elf
    \end{bashcode}
    \bigbreak \noindent 
    At the GDB prompt, enter:
    \begin{bashcode}
        target extended-remote localhost:3333
        monitor reset halt
        load
        monitor reset halt
        break main
        continue
    \end{bashcode}
    \begin{itemize}
        \item \textbf{Target:}  is a command used to specify the execution environment of the program you are debugging. It connects your local GDB session (the host) to a process, core file, or a separate machine (the target). 
        \item \textbf{load}:  The load command in GDB is primarily used during remote debugging to transfer and upload an executable file directly onto a remote target system (such as an embedded microcontroller or a remote board)
        \item \textbf{monitor reset halt}: In GDB, the command monitor reset halt is used during embedded firmware debugging to reset the target hardware and immediately halt the CPU at the very first instruction
            \bigbreak \noindent 
            Monitor tells GDB to pass the trailing command directly to the underlying hardware debug server (such as OpenOCD or pyOCD) rather than executing it within GDB itself
            \bigbreak \noindent 
            Reset halt instructs the debug server to assert the hardware reset line and stop the processor core immediately upon release. This prevents the MCU from executing any application code or bootloaders, keeping it in a clean, predictable state
            \bigbreak \noindent 
            The first resets and halts the MCU so OpenOCD can program Flash while the processor is not executing code.
            \bigbreak \noindent 
            The second resets the MCU after programming so the newly loaded firmware begins with a clean hardware state:
            \begin{itemize}
                \item Stack pointer is reloaded.
                \item Program counter returns to Reset\_Handler.
                \item Global initialization runs again.
                \item Peripherals return to their reset state.
            \end{itemize}
            The first reset is not always strictly necessary because OpenOCD can normally halt the target automatically during load.
    \end{itemize}
    The complete connection is
    \begin{center}
        GDB $\to$ OpenOCD $\to$ onboard ST-LINK $\to$ SWD $\to$ STM32F446RE
    \end{center}
    build/intro.elf contains:
    \begin{itemize}
        \item ARM machine instructions.
        \item The addresses where sections belong in Flash and SRAM.
        \item The interrupt vector table.
        \item Initialized global data.
        \item Function and variable names.
        \item Source-line debugging information.
        \item Section metadata such as .text, .data, .bss, and .isr_vector.
    \end{itemize}
    The MCU never parses the ELF file. OpenOCD reads the ELF metadata, extracts its loadable machine-code sections, and writes them into the correct Flash locations. After reset, the Cortex-M4 fetches instructions directly from Flash.
    \bigbreak \noindent 
    You can inspect the ELF structure with:
    \bigbreak \noindent 
    \begin{bashcode}
        arm-none-eabi-readelf -h build/*.elf
        arm-none-eabi-objdump -h build/*.elf
        arm-none-eabi-nm build/*.elf
    \end{bashcode}
    \bigbreak \noindent 
    Once the firmware is flashed, the STM32 runs independently. GDB and OpenOCD do not need to remain connected. With
    \bigbreak \noindent 
    \begin{bashcode}
        openocd \
            -f board/st_nucleo_f4.cfg \
            -c "program build/intro.elf verify reset exit"
    \end{bashcode}
    \bigbreak \noindent 
    The final reset exit means:
    \begin{itemize}
        \item Reset the STM32.
        \item Allow it to run.
        \item Exit OpenOCD.
    \end{itemize}
    After that, the program remains in Flash and automatically starts whenever you:
    \begin{itemize}
        \item Press the board’s reset button.
        \item Disconnect and reconnect power.
        \item Turn the board off and back on.
    \end{itemize}

















\end{document}
